% META: {"id": "TSPE-SUITLIM-EV-B", "chapitre": "TSPE-SUITES-LIMITES", "type_objet": "evaluation", "version": "B", "duree_min": 55, "bareme_total": 20, "capacites_codes": ["C1", "C2", "C3", "C4", "C5", "C6", "C7"], "competences": ["chercher", "modeliser", "representer", "calculer", "raisonner", "communiquer"], "mode_creation": "ex_nihilo", "sources_inspiration": [], "status": "generated"}
% BEGIN-VERIFY
% from sympy import *
% import math
% n = symbols('n', integer=True, nonnegative=True)
%
% # ---- Exercice 1 (5 pts) : Limites (C1, C5) ----
% u1 = (4*n**2 - n)/(2*n**2 + 5)
% assert limit(u1, n, oo) == 2
% u2 = 3*Rational(2,5)**n + 7
% assert limit(u2, n, oo) == 7
% u3 = 2*n**2 - 5*n + 3
% assert limit(u3, n, oo) == oo
%
% # ---- Exercice 2 (5 pts) : Recurrence (C2) ----
% for k in range(15):
%     assert 3**k >= 2*k + 1
%
% # ---- Exercice 3 (5 pts) : Modelisation (C3) ----
% a3 = Rational(4, 5)
% b3 = 60
% l3 = b3 / (1 - a3)
% assert l3 == 300
% u3_expr = 300 - 100*a3**n
% assert u3_expr.subs(n, 0) == 200
% assert u3_expr.subs(n, 1) == 220
% assert limit(u3_expr, n, oo) == 300
%
% # ---- Exercice 4 (5 pts) : Suite recurrente (C4, C6) ----
% u4 = [2.0]
% for i in range(20):
%     u4.append(math.sqrt(2*u4[-1] + 3))
% assert abs(u4[-1] - 3) < 1e-6
% END-VERIFY

%---------------------------------------------------------------
% Duree : 55 minutes — Bareme : 20 points
% Toute reponse doit etre justifiee ; les resultats seuls ne sont pas acceptes.
%---------------------------------------------------------------

\begin{evaluation}{TSPE-SUITLIM-EV-B}{55}

%===============================================================
% EXERCICE 1 — Limites de suites (5 points)
%===============================================================

\begin{exercice}{TSPE-SUITLIM-EV-B-EX1}{2}{12}

\textbf{Exercice 1} (5 points) — \textit{Capacites C1, C5}

\medskip

Determiner les limites des suites suivantes en justifiant.

\begin{enumerate}
  \item (1,5 pt) $u_n = \dfrac{4n^2 - n}{2n^2 + 5}$
  \item (1,5 pt) $v_n = 3 \times \left(\dfrac{2}{5}\right)^n + 7$
  \item (2 pts) $w_n = 2n^2 - 5n + 3$ (montrer que $w_n \geq n^2$ pour $n$ assez grand)
\end{enumerate}

\end{exercice}

%===============================================================
% EXERCICE 2 — Recurrence (5 points)
%===============================================================

\begin{exercice}{TSPE-SUITLIM-EV-B-EX2}{2}{12}

\textbf{Exercice 2} (5 points) — \textit{Capacite C2}

\medskip

\begin{enumerate}
  \item (3 pts) Demontrer par recurrence que pour tout entier naturel $n$, $3^n \geq 2n + 1$.
  \item (2 pts) En deduire $\lim_{n \to +\infty} 3^n$.
\end{enumerate}

\end{exercice}

%===============================================================
% EXERCICE 3 — Modelisation (5 points)
%===============================================================

\begin{exercice}{TSPE-SUITLIM-EV-B-EX3}{2}{15}

\textbf{Exercice 3} (5 points) — \textit{Capacite C3}

\medskip

Un bassin contient $200$ litres d'eau. Chaque jour, $20\,\%$ de l'eau s'evapore et on ajoute $60$~litres. On note $u_n$ le volume au debut du jour $n$.

\begin{enumerate}
  \item (1 pt) Justifier que $u_0 = 200$ et $u_{n+1} = 0{,}8\,u_n + 60$.
  \item (1 pt) Calculer $u_1$ et $u_2$.
  \item (2 pts) Poser $v_n = u_n - \ell$ avec $\ell$ adequat. Exprimer $u_n$.
  \item (1 pt) Determiner $\lim u_n$ et interpreter.
\end{enumerate}

\end{exercice}

%===============================================================
% EXERCICE 4 — Suite recurrente (5 points)
%===============================================================

\begin{exercice}{TSPE-SUITLIM-EV-B-EX4}{3}{16}

\textbf{Exercice 4} (5 points) — \textit{Capacites C4, C6}

\medskip

Soit $(u_n)$ definie par $u_0 = 2$ et $u_{n+1} = \sqrt{2u_n + 3}$.

\begin{enumerate}
  \item (1 pt) Calculer $u_1$, $u_2$.
  \item (1,5 pt) Montrer par recurrence que $0 \leq u_n \leq 3$.
  \item (1,5 pt) Montrer que $(u_n)$ est croissante.
  \item (1 pt) En deduire que $(u_n)$ converge et determiner sa limite.
\end{enumerate}

\end{exercice}

\end{evaluation}
